\documentclass[literature.tex]{subfiles}

\begin{document}
\drama{Romeo a Julie}{William Shakespeare}
\teorieA{
drama, tragédie (tragédie osudu s prvky romantické lásky a pomsty)
}{
Anglická renesance (elizabethanská doba), humanismus.
Shakespeare kombinuje prvky středověké tragédie, petrarkovské lásky a italské novely.
}{
Blankvers (pětistopý jambický verš) u vyšších postav, proza u nižších vrstev (chůva).
Bohatý poetický jazyk s metaforami, přirovnáními, slovními hříčkami a citově zabarvenými výrazy.
Kontrast vysokého stylu (Romeo, Julie) a lidové mluvy.
}{
Řečnické otázky (\enquote{Romeo, proč's Romeo?}), metafora, přirovnání, oxymóron (sladká bolest lásky), dramatická ironie.
}{
Symbolismus (balkon = hranice mezi rody, jed = osud), alegorie (feud = nenávist společnosti), foreshadowing (prolog sonet předpovídá tragický konec).
}{
Pět dějství s chronologickým postupem.
Směs tragického a komického (první akty mají komediální prvky).
Rychlé dialogy, monology (balkonová scéna) a sbor (prolog).
Tragikomický tón – láska a smrt jsou propojeny.
}{
\wemph{Jul.} Ach, žel! \\
\wemph{Rom.} Teď mluví! -- Promluv poznovu, \par
ó jasný anděle! -- neb nad hlavou \par
mi vystupuješ zářná z noční tmy, \par
jak okřídlený posel nebeský \par
v sloup obráceným očím žasnoucích \par
smrtelných lidí, padajících na znak, \par
by dívali se naň, jak osedlal \par
si lenivě plynoucí oblaky \par
a po nebeských ňadrech vesluje. \\
\wemph{Jul.} Romeo, Ó Romeo! -- proč's Romeo? \par
Své jméno zapři, otce zřekni se, \par
neb, nechceš-li, mně lásku přísahej, \par
a nechci dál být Capuletova.
}
\teorieB{
\wtem{Romeo}{mladý, romantický, impulzivní, zamilovaný do lásky samotné (nejprve Rosalina, pak Julie).
Statečný, ale ukvapený.}
\wtem{Julie}{mladá (téměř 14 let), inteligentní, odvážná, tvrdohlavá.
Zpočátku poslušná, později se vzepře rodičům kvůli lásce.}
\wtem{Tybalt}{Juliin bratranec, horkokrevný, bojovný, symbol rodinné nenávisti.}
\wtem{Merkucio}{Romeův přítel, vtipný, cynický, bojovný.
Jeho smrt urychlí tragédii.}
\wtem{BratrVavřinec}{františkán, moudrý, však naivní.
Snaží se pomoci, jeho plán selže.}
\wtem{Chůva}{Juliina věrná služebnice, komická, lidová postava, praktická.}
}{
Ve Veroně žijí dva znepřátelené rody -- Montekové a Kapuletové.
Romeo (Montek) se na plese zamiluje do Julie (Kapuletové).
Oba se tajně oddají u bratra Vavřince.
Po pouliční bitce Romeo zabije Tybalta (Juliina bratrance) a je vyhoštěn.
Julie má být provdána za Parise.
Bratr Vavřinec jí dá uspávací nápoj, aby vypadala mrtvá.
Plán selže -- Romeo se dozví o „smrti“ Julie, vrátí se, zabije Parise a spáchá sebevraždu jedem.
Julie se probudí a probodne se Romeovou dýkou.
Obě rodiny se nad mrtvými těly usmíří.
}{
Pět dějství s prologem (sonet, který shrnuje děj).
Chronologický postup (události proběhnou během několika dní).
Prostor: Verona (Itálie).
Čas: 14.–16. století (elizabethanská doba).
}{
Tragédie nešťastné lásky, kterou ničí rodinná nenávist.
Síla lásky překonává společenské bariéry, ale osud („star-crossed lovers“) a ukvapenost vedou k tragédii.
Kritika slepé nenávisti a rodičovské autority.
Láska jako osvobozující, ale destruktivní síla.
Závěrečné usmíření ukazuje, že smrt mladých přináší mír.
}
\historie{
Alžbětinská Anglie (vláda Alžběty I., konec 16. století) -- náboženské napětí (katolíci vs. protestanti), stabilita po občanských válkách.
Shakespeare kritizuje feudální spory skrz italské prostředí.
}{}{
Inspirováno italskými novelami (Matteo Bandello) a anglickou básní Arthura Brooka (1562).
Popularizovalo téma zakázané lásky.
}{}{
Raná Shakespearova tragédie.
Kombinuje komediální prvky (první akty) s tragédií.
Ovlivnil pozdější romantismus a moderní adaptace (West Side Story, filmy).
\newpage}{
William Shakespeare (1564 Stratford-upon-Avon – 1616).
Nejvýznamnější anglický dramatik.
Psal v alžbětinské době. Autor 37 her, 154 sonetů.
Jeho tvorba se dělí na rané (komedie, historické hry), střední (tragédie jako Hamlet) a pozdní období.
}{
Mezi nejznámější díla patří \textit{Hamlet}, \textit{Othello}, \textit{Sen noci svatojánské}, \textit{Král Lear} a \textit{Makbeth}.
}
\kritika{}{}{
Nadčasové – zakázaná láska, generační konflikt, nenávist mezi skupinami.
Adaptace do moderních kontextů (rasismus, gangy) ukazují, že téma zůstává relevantní.
}
\nazor{
Romeo a Julie je klasika, která stále zasáhne.
Balkonová scéna a intenzita mladé lásky jsou nádherné, ale nejsilnější je tragédie způsobená rodinnou nenávistí a ukvapeností.
Shakespeare mistrně ukazuje, jak láska může být osvobozující i destruktivní.
Prolog sonet hned prozradí konec, přesto napětí funguje.
Je to varování před slepou nenávistí a oslavou lásky, která překonává společenské bariéry.
I po stovkách let patří k nejsilnějším milostným příběhům.
}
\explain{
\textbl{Tragédie --} dramatický žánr s tragickým koncem hrdinů a morálním ponaučením;
\textbl{Blankvers --} nerýmovaný pětistopý jambický verš;
\textbl{Foreshadowing --} předznamenání budoucích událostí (prolog).
}
\wpage
\end{document}